% GENERATED FILE. Edit canonical lessons, then run npm run generate:books.
% canonical-source-hash: fnv1a64:fe900fedc9d86adf
% canonical-lessons: KA-C341-vamsa, KA-C341-purvaja, KA-C341-talemaru, KA-C341-manetana, KA-C341-arasa

\chapter{Our Family Line}
\label{ch:ka-our-family-line}
\hlchaptermodality{341}

\begin{chapteropening}
\textbf{By the end of this chapter:} \emph{I can say a lineage, an ancestor, a family generation, a family line and a monarch.}

Complete the last lesson of chapter 341: I can say a lineage, an ancestor, a family generation, a family line and a monarch.
\end{chapteropening}

\section[\texorpdfstring{\kn{ವಂಶ} (vaṃśa)}{ವಂಶ (vaṃśa)}]{\kn{ವಂಶ} (vaṃśa) — a lineage; a dynasty}

\label{lesson:KA-C341-vamsa}



\begin{quote}
Before the new one: say the Kannada for a demon, then the Kannada for a curse.
\end{quote}

\subsection*{You'll want to know: \kn{ವಂಶ}}
\textbf{\kn{ವಂಶ}} — \emph{vaṃśa} — \textquotedblleft{}a lineage; a dynasty\textquotedblright{}.

\begin{quote}
\kn{ಆ} \kn{ವಂಶದ} \kn{ರಾಜರು} \kn{ಮುನ್ನೂರು} \kn{ವರ್ಷ} \kn{ಆಳಿದರು}.
\end{quote}

\begin{quote}
\emph{ā vaṃśada rājaru munnūru varṣa āḷidaru.}
\end{quote}

\begin{quote}
The kings of that dynasty ruled for three hundred years.
\end{quote}

\kn{ವಂಶವೃಕ್ಷ} is a family tree.

\begin{cousinweb}
\textbf{\kn{ವಂಶ}} (\emph{vaṃśa}) is borrowed from Sanskrit \emph{vaṃśa}, \textquotedblleft{}a bamboo cane\textquotedblright{}, and from its joints, \textquotedblleft{}a line of descent, a race\textquotedblright{} (Monier-Williams).
\end{cousinweb}

\begin{grammarlens}[title={What you've built}]
One more word for telling what happened.
\end{grammarlens}

\subsection*{Your turn}
\begin{itemize}
\raggedright
  \item \emph{Say it:} \emph{vaṃśa}
  \item \emph{Say it:} \emph{vaṃśa}, once more
  \item \emph{Say it:} \emph{ā vaṃśada rājaru munnūru varṣa āḷidaru.} — The kings of that dynasty ruled for three hundred years.
  \item \emph{Recall:} say the Kannada for a demon, then the Kannada for a curse, then say \emph{vaṃśa} again
\end{itemize}

\begin{tcolorbox}[breakable,colback=teal!4,colframe=teal!35!black,title={Before you move on}]
What does \kn{ವಂಶ} mean? (\textquotedblleft{}A lineage; a dynasty\textquotedblright{}.) Say it once more.
\end{tcolorbox}
\section[\texorpdfstring{\kn{ಪೂರ್ವಜ} (pūrvaja)}{ಪೂರ್ವಜ (pūrvaja)}]{\kn{ಪೂರ್ವಜ} (pūrvaja) — an ancestor}

\label{lesson:KA-C341-purvaja}



\begin{quote}
Before the new one: say the Kannada for a curse, then the Kannada for a lineage.
\end{quote}

\subsection*{You'll want to know: \kn{ಪೂರ್ವಜ}}
\textbf{\kn{ಪೂರ್ವಜ}} — \emph{pūrvaja} — \textquotedblleft{}an ancestor\textquotedblright{}.

\begin{quote}
\kn{ನಮ್ಮ} \kn{ಪೂರ್ವಜರು} \kn{ಈ} \kn{ಊರಿಗೆ} \kn{ಬಂದು} \kn{ಮನೆ} \kn{ಕಟ್ಟಿದರು}.
\end{quote}

\begin{quote}
\emph{namma pūrvajaru ī ūrige bandu mane kaṭṭidaru.}
\end{quote}

\begin{quote}
Our ancestors came to this village and built houses.
\end{quote}

Plural: \kn{ಪೂರ್ವಜರು}, \emph{pūrvajaru}.

\kn{ಪೂರ್ವ}, before, and \kn{ಜ}, born: one born before.

\begin{cousinweb}
\textbf{\kn{ಪೂರ್ವಜ}} (\emph{pūrvaja}) is borrowed from Sanskrit \emph{pūrvaja}, \textquotedblleft{}born before, former, ancient\textquotedblright{} (Monier-Williams).
\end{cousinweb}

\begin{grammarlens}[title={What you've built}]
One more word for telling what happened.
\end{grammarlens}

\subsection*{Your turn}
\begin{itemize}
\raggedright
  \item \emph{Say it:} \emph{pūrvaja}
  \item \emph{Say it:} \emph{pūrvaja}, once more
  \item \emph{Say it:} \emph{namma pūrvajaru ī ūrige bandu mane kaṭṭidaru.} — Our ancestors came to this village and built houses.
  \item \emph{Recall:} say the Kannada for a curse, then the Kannada for a lineage, then say \emph{pūrvaja} again
\end{itemize}

\begin{tcolorbox}[breakable,colback=teal!4,colframe=teal!35!black,title={Before you move on}]
What does \kn{ಪೂರ್ವಜ} mean? (\textquotedblleft{}An ancestor\textquotedblright{}.) Say it once more.
\end{tcolorbox}
\section[\texorpdfstring{\kn{ತಲೆಮಾರು} (talemāru)}{ತಲೆಮಾರು (talemāru)}]{\kn{ತಲೆಮಾರು} (talemāru) — a family generation}

\label{lesson:KA-C341-talemaru}



\begin{quote}
Before the new one: say the Kannada for a lineage, then the Kannada for an ancestor.
\end{quote}

\subsection*{You'll want to know: \kn{ತಲೆಮಾರು}}
\textbf{\kn{ತಲೆಮಾರು}} — \emph{talemāru} — \textquotedblleft{}a family generation\textquotedblright{}.

\begin{quote}
\kn{ಈ} \kn{ಮನೆಯಲ್ಲಿ} \kn{ನಾಲ್ಕು} \kn{ತಲೆಮಾರಿನವರು} \kn{ಒಟ್ಟಿಗೆ} \kn{ಇದ್ದರು}.
\end{quote}

\begin{quote}
\emph{ī maneyalli nālku talemārinavaru oṭṭige iddaru.}
\end{quote}

\begin{quote}
Four generations lived together in this house.
\end{quote}

Plural: \kn{ತಲೆಮಾರುಗಳು}, \emph{talemārugaḷu}.

Grandparents, parents and children are three \kn{ತಲೆಮಾರು}.

\begin{cousinweb}
\textbf{\kn{ತಲೆಮಾರು}} (\emph{talemāru}) is built on native \emph{tale}, \textquotedblleft{}head, top\textquotedblright{} (DEDR 3103); the origin of the second part, -\kn{ಮಾರು}, is not settled.
\end{cousinweb}

\begin{grammarlens}[title={What you've built}]
One more word for telling what happened.
\end{grammarlens}

\subsection*{Your turn}
\begin{itemize}
\raggedright
  \item \emph{Say it:} \emph{talemāru}
  \item \emph{Say it:} \emph{talemāru}, once more
  \item \emph{Say it:} \emph{ī maneyalli nālku talemārinavaru oṭṭige iddaru.} — Four generations lived together in this house.
  \item \emph{Recall:} say the Kannada for a lineage, then the Kannada for an ancestor, then say \emph{talemāru} again
\end{itemize}

\begin{tcolorbox}[breakable,colback=teal!4,colframe=teal!35!black,title={Before you move on}]
What does \kn{ತಲೆಮಾರು} mean? (\textquotedblleft{}A family generation\textquotedblright{}.) Say it once more.
\end{tcolorbox}
\section[\texorpdfstring{\kn{ಮನೆತನ} (manetana)}{ಮನೆತನ (manetana)}]{\kn{ಮನೆತನ} (manetana) — a family line; a household}

\label{lesson:KA-C341-manetana}



\begin{quote}
Before the new one: say the Kannada for an ancestor, then the Kannada for a family generation.
\end{quote}

\subsection*{You'll want to know: \kn{ಮನೆತನ}}
\textbf{\kn{ಮನೆತನ}} — \emph{manetana} — \textquotedblleft{}a family line; a household\textquotedblright{}.

\begin{quote}
\kn{ಅವರದು} \kn{ಹಳೆಯ} \kn{ಶ್ರೀಮಂತ} \kn{ಮನೆತನ}.
\end{quote}

\begin{quote}
\emph{avaradu haḷeya śrīmanta manetana.}
\end{quote}

\begin{quote}
Theirs is an old wealthy family.
\end{quote}

\kn{ಮನೆ}, house, with -\kn{ತನ}, \textquotedblleft{}-ness\textquotedblright{}: the house as a family through time.

\begin{cousinweb}
\textbf{\kn{ಮನೆತನ}} (\emph{manetana}) is native Dravidian: DEDR 4776 lists Kannada \emph{manetana}, \textquotedblleft{}household, household life\textquotedblright{}, built on \emph{mane}, \textquotedblleft{}a house\textquotedblright{}.
\end{cousinweb}

\begin{grammarlens}[title={What you've built}]
One more word for telling what happened.
\end{grammarlens}

\subsection*{Your turn}
\begin{itemize}
\raggedright
  \item \emph{Say it:} \emph{manetana}
  \item \emph{Say it:} \emph{manetana}, once more
  \item \emph{Say it:} \emph{avaradu haḷeya śrīmanta manetana.} — Theirs is an old wealthy family.
  \item \emph{Recall:} say the Kannada for an ancestor, then the Kannada for a family generation, then say \emph{manetana} again
\end{itemize}

\begin{tcolorbox}[breakable,colback=teal!4,colframe=teal!35!black,title={Before you move on}]
What does \kn{ಮನೆತನ} mean? (\textquotedblleft{}A family line; a household\textquotedblright{}.) Say it once more.
\end{tcolorbox}
\section[\texorpdfstring{\kn{ಅರಸ} (arasa)}{ಅರಸ (arasa)}]{\kn{ಅರಸ} (arasa) — a monarch; a king}

\label{lesson:KA-C341-arasa}



\begin{quote}
Before the new one: say the Kannada for a family generation, then the Kannada for a family line.
\end{quote}

\subsection*{You'll want to know: \kn{ಅರಸ}}
\textbf{\kn{ಅರಸ}} — \emph{arasa} — \textquotedblleft{}a monarch; a king\textquotedblright{}.

\begin{quote}
\kn{ಆ} \kn{ಅರಸ} \kn{ಪ್ರಜೆಗಳಿಗಾಗಿ} \kn{ಕೆರೆಗಳನ್ನು} \kn{ಕಟ್ಟಿಸಿದನು}.
\end{quote}

\begin{quote}
\emph{ā arasa prajegaḷigāgi keregaḷannu kaṭṭisidanu.}
\end{quote}

\begin{quote}
That king had lakes built for his people.
\end{quote}

Plural: \kn{ಅರಸರು}, \emph{arasaru}.

An older Kannada form beside \kn{ರಾಜ}; \kn{ಅರಮನೆ}, a palace, is the king's house.

\begin{cousinweb}
\textbf{\kn{ಅರಸ}} (\emph{arasa}) goes back to Sanskrit \emph{rājan}, \textquotedblleft{}a chieftain, a king\textquotedblright{} (CDIAL 10679): DEDR 201 lists Kannada \emph{arasa}, \emph{arasu}, \textquotedblleft{}a king, a lord\textquotedblright{}, and \emph{ara-mane}, \textquotedblleft{}a palace\textquotedblright{}, and derives them from the Sanskrit word. It is an old reshaped loan, beside the learned \kn{ರಾಜ}.
\end{cousinweb}

\begin{grammarlens}[title={What you've built}]
One more word for telling what happened.
\end{grammarlens}

\subsection*{Your turn}
\begin{itemize}
\raggedright
  \item \emph{Say it:} \emph{arasa}
  \item \emph{Say it:} \emph{arasa}, once more
  \item \emph{Say it:} \emph{ā arasa prajegaḷigāgi keregaḷannu kaṭṭisidanu.} — That king had lakes built for his people.
  \item \emph{Recall:} say the Kannada for a family generation, then the Kannada for a family line, then say \emph{arasa} again
\end{itemize}

\begin{tcolorbox}[breakable,colback=teal!4,colframe=teal!35!black,title={Before you move on}]
What does \kn{ಅರಸ} mean? (\textquotedblleft{}A monarch; a king\textquotedblright{}.) Say it once more.
\end{tcolorbox}
